\section*{अध्याय \ref{ch:b2:reduction} --- अंतःरूपांतरणों का समानयन}

\begin{solution}{exo:b2:reduction:1}
$A$: $\chi_A = X^2 - 2X - 3 = (X - 3)(X + 1)$। \omterm{def:b2:reduction:eigen}{अभिलक्षणिक सदिश}: $3$ के लिए $(1,1)$; $-1$ के लिए $(1,-1)$। अतः $P = \begin{pmatrix} 1 & 1\\ 1
& -1\end{pmatrix}$ से
$P^{-1}AP = \operatorname{diag}(3, -1)$ मिलता है।

$B = J - I$, जहाँ $J$ सर्व-एक आव्यूह है। $J$ की कोटि $1$ है, जिसमें $v = (1,1,1)$ के लिए
$Jv = 3v$ और समतल $x + y + z =
0$ पर $Jw = 0$: अतः $B$ का \omterm{def:b2:reduction:eigen}{वर्णक्रम} $\{2, -1\}$ है, जिसकी अभिलक्षणिक
समष्टियाँ $\operatorname{Vect}(1,1,1)$ (विमा $1$) और $\{x + y + z = 0\}$ (विमा $2$, आधार $(1,-1,0), (1,0,-1)$) हैं। इन
तीन स्तंभों वाला $P$ $P^{-1}BP = \operatorname{diag}(2, -1, -1)$ देता है।
\end{solution}

\begin{solution}{exo:b2:reduction:2}
\emph{अभिलक्षणिक समष्टियाँ:} $\chi_C = (X-1)^2$, एकमात्र \omterm{def:b2:reduction:eigen}{अभिलक्षणिक मान} $1$; और $\ker(C - I) = \ker\begin{pmatrix} 0&1\\ 0&0\end{pmatrix}$
रेखा $\operatorname{Vect}(e_1)$ है: विमा $1 < 2 = m_1$, अतः \omterm{def:b2:reduction:diag}{विकर्णनीय} नहीं (\cref{thm:b2:reduction:diagcrit})।

\emph{\omterm{def:b2:reduction:polyu}{न्यूनतम बहुपद}:} $\mu_C$ $(X-1)^2$ को विभाजित करता है और $C \neq I$ को भी, अतः
$\mu_C = (X-1)^2$: दोहरा मूल, इसलिए \omterm{def:b2:reduction:diag}{विकर्णनीय} नहीं (\cref{cor:b2:reduction:minpolycrit})।
\end{solution}

\begin{solution}{exo:b2:reduction:3}
$X^2 - 5X + 6 = (X-2)(X-3)$: सरल मूलों के साथ पूरी तरह गुणनखंडित, अतः $u$ \omterm{def:b2:reduction:diag}{विकर्णनीय} है (\cref{cor:b2:reduction:minpolycrit}),
और $\operatorname{Sp}(u) \subseteq \{2, 3\}$। संभव \omterm{def:b2:reduction:eigen}{वर्णक्रम}: $\{2\}$ ($u = 2\,\mathrm{id}$), $\{3\}$ ($u = 3\,\mathrm{id}$), अथवा $\{2, 3\}$।

घातें: $u^k = a_k\,\mathrm{id} + b_k\,u$ खोजिए। अभिलक्षणिक समष्टियों पर यह $2^k = a_k + 2b_k$ और $3^k = a_k + 3b_k$ पढ़ा जाता है:
हल करने पर $b_k
= 3^k - 2^k$, $a_k = 3\cdot2^k - 2\cdot 3^k$:
\[
u^k = (3\cdot 2^k - 2\cdot 3^k)\,\mathrm{id} + (3^k - 2^k)\, u .
\]
(तीनों वर्णक्रमों के लिए मान्य: सर्वसमिकाएँ अभिलक्षणिक-मान-दर-अभिलक्षणिक-मान
सत्य हैं।)
\end{solution}

\begin{solution}{exo:b2:reduction:4}
$u$ \omterm{def:b2:reduction:diag}{विकर्णनीय}: \omterm{def:b2:reduction:eigen}{वर्णक्रम} पर $P = \prod_{\lambda}(X - \lambda)$ $u$ को विलुप्त करता है, पूरी तरह
गुणनखंडित है, और उसके मूल सरल हैं। तब $P(u|_F) =
P(u)|_F = 0$: प्रतिबंधन भी सरल मूलों वाले
पूरी तरह गुणनखंडित बहुपद से विलुप्त होता है, अतः \omterm{def:b2:reduction:diag}{विकर्णनीय} है (\cref{cor:b2:reduction:minpolycrit})।
\end{solution}

\begin{solution}{exo:b2:reduction:5}
$\chi_A = X^2 - X - 1$, जिसके मूल $\varphi$ और $\psi$ हैं (भिन्न): अतः \omterm{def:b2:reduction:diag}{विकर्णनीय}, और \omterm{def:b2:reduction:eigen}{अभिलक्षणिक
सदिश} $(\varphi, 1)$ तथा $(\psi, 1)$। पुनरावृत्ति $\begin{pmatrix} F_{n+1}\\ F_n \end{pmatrix} =
A^n \begin{pmatrix}1\\ 0\end{pmatrix}$ देती है। $(1, 0)$ को अभिलक्षणिक सदिशों
पर अपघटित कीजिए: $\varphi - \psi = \sqrt5$ सहित $(1,0) = \frac{1}{\varphi - \psi}\bigl((\varphi, 1) -
(\psi, 1)\bigr)$। $A^n$ लगाने से हर अभिलक्षणिक घटक अपने
\omterm{def:b2:reduction:eigen}{अभिलक्षणिक मान} की $n$-वीं घात से गुणा हो जाता है; दूसरा निर्देशांक
पढ़ने पर:
\[
F_n = \frac{\varphi^n - \psi^n}{\sqrt 5} .
\]
(जाँच: $n = 1$ से $\frac{\varphi - \psi}{\sqrt5} = 1$ मिलता है।)
\end{solution}

\begin{solution}{exo:b2:reduction:6}
मान लीजिए $P = \prod_i (X - \mu_i)$ $u^2$ को विलुप्त करता है और सरल मूलों $\mu_i$ (अर्थात् $u^2$ का
\omterm{def:b2:reduction:eigen}{वर्णक्रम}) के साथ पूरी तरह गुणनखंडित होता है। चूँकि $u$ प्रतिलोमनीय है,
$0$ $u^2$ का \omterm{def:b2:reduction:eigen}{अभिलक्षणिक मान} नहीं है ($\det u^2 = (\det u)^2 \neq 0$), अतः सभी $\mu_i \neq 0$। तब
\[
Q(X) = \prod_i (X^2 - \mu_i) = \prod_i (X - \sqrt{\mu_i})(X +
\sqrt{\mu_i})
\]
$u$ को विलुप्त करता है: $\;Q(u) = \prod_i (u^2 - \mu_i\,\mathrm{id}) =
P(u^2) = 0$। उसके मूल $\pm
\sqrt{\mu_i}$ (सम्मिश्र वर्गमूल) परस्पर
भिन्न हैं क्योंकि $\mu_i$ भिन्न और अशून्य हैं ($\sqrt{\mu_i} = -\sqrt{\mu_j}$ से $\mu_i = \mu_j$ मिल जाता)।
पूरी तरह गुणनखंडित और सरल मूल: अतः $u$ \omterm{def:b2:reduction:diag}{विकर्णनीय} है।

प्रतिलोमनीयता के बिना प्रतिउदाहरण: $u = \begin{pmatrix} 0 & 1\\
0 & 0\end{pmatrix}$: $u^2 = 0$ \omterm{def:b2:reduction:diag}{विकर्णनीय} है, $u$ नहीं।
\end{solution}

\begin{solution}{exo:b2:reduction:7}
$A = 2I + N$, जहाँ $N$ स्थानांतरण है ($N e_2 = e_1$, $Ne_3 = e_2$), $N^3 = 0$, $N^2 = E_{13}$: यह \emph{ही} \omterm{thm:b2:reduction:dunford}{डनफ़र्ड
अपघटन} है ($2I$ विकर्ण, $N$ शून्यंभावी, और वे क्रमविनिमेय हैं;
अद्वितीयता से यही वह अपघटन है)। क्रमविनिमेय पदों के साथ द्विपद:
\[
A^k = 2^k I + k 2^{k-1} N + \binom k2 2^{k-2} N^2
= \begin{pmatrix}
2^k & k2^{k-1} & \binom k2 2^{k-2}\\
0 & 2^k & k2^{k-1}\\
0 & 0 & 2^k
\end{pmatrix},
\]
\[
\eu^{tA} = \eu^{2t}\Bigl(I + tN + \frac{t^2}{2}N^2\Bigr)
= \eu^{2t}\begin{pmatrix}
1 & t & t^2/2\\
0 & 1 & t\\
0 & 0 & 1
\end{pmatrix}.
\]
\end{solution}

\begin{solution}{exo:b2:reduction:8}
$X^k - 1$ $A$ को विलुप्त करता है और $\C$ पर $k$ भिन्न मूलों $\eu^{2\iu\pi j/k}$ के साथ
गुणनखंडित होता है: अतः $A$ \omterm{def:b2:reduction:diag}{विकर्णनीय} है (\cref{cor:b2:reduction:minpolycrit}) और उसके \omterm{def:b2:reduction:eigen}{अभिलक्षणिक मान},
जो $X^k - 1$ के मूल हैं, इकाई के $k$-वें मूल हैं।

यदि इसके अतिरिक्त $A$ इकाई विकर्ण वाले किसी त्रिभुजाकार आव्यूह के सदृश
है: तो सभी \omterm{def:b2:reduction:eigen}{अभिलक्षणिक मान} $1$ के बराबर हैं, और $A$, जो एकमात्र
\omterm{def:b2:reduction:eigen}{अभिलक्षणिक मान} $1$ के साथ \omterm{def:b2:reduction:diag}{विकर्णनीय} है, $P\,I\,P^{-1} = I$ है।
\end{solution}

\begin{solution}{exo:b2:reduction:9}
$E = \bigoplus_\lambda E_\lambda(u)$ लिखिए (\cref{thm:b2:reduction:diagcrit})। प्रत्येक $E_\lambda(u)$ $v$-स्थायी है: $x \in E_\lambda$ के लिए $u(v(x)) = v(u(x)) = \lambda
v(x)$। $v$ का
$E_\lambda(u)$ पर प्रतिबंधन \omterm{def:b2:reduction:diag}{विकर्णनीय} है (\cref{exo:b2:reduction:4}): $E_\lambda(u)$ का ऐसा आधार चुनिए जो
$v$-अभिलक्षणिक सदिशों से बना हो। सभी $\lambda$ पर इन आधारों को जोड़ देने से
$E$ का ऐसा आधार मिलता है जिसके सदिश \emph{दोनों} $u$ ($E_\lambda(u)$ में होने
से) और $v$ (रचना से) के \omterm{def:b2:reduction:eigen}{अभिलक्षणिक सदिश} हैं।
\end{solution}

\begin{solution}{exo:b2:reduction:10}
($\Rightarrow$) मान लीजिए $u$ \omterm{def:b2:reduction:diag}{विकर्णनीय} है और $F$ स्थायी। तब $u|_F$ \omterm{def:b2:reduction:diag}{विकर्णनीय} है
(\cref{exo:b2:reduction:4}): $F$ का कोई आधार अभिलक्षणिक सदिशों से बना है, जो प्रत्येक वैश्विक
\omterm{def:b2:reduction:eigen}{अभिलक्षणिक समष्टि} $E_\lambda$ के भीतर $E_\lambda$ के आधार तक बढ़ाया जा सकता है ($E_\lambda$ के
भीतर अपूर्ण आधार प्रमेय, वहाँ पड़े $F$ के आधार के भाग से आरंभ करके ---
ध्यान दीजिए कि $u|_F$ के \omterm{def:b2:reduction:diag}{विकर्णनीय} होने से $F = \bigoplus_\lambda (F \cap E_\lambda)$)। जोड़े गए सदिश एक स्थायी
संपूरक फैलाते हैं (हर सदिश किसी न किसी $E_\lambda$ में है, अतः उनका विस्तार
$u$-स्थायी है)।

($\Leftarrow$) मान लीजिए $F = \sum_\lambda E_\lambda(u)$ (एक स्थायी उपसमष्टि) और $G$ उसका स्थायी संपूरक है।
यदि $G \neq \{0\}$: तो $\chi_{u|_G}$ $\C$ पर गुणनखंडित होता है, अतः $u|_G$ का कोई \omterm{def:b2:reduction:eigen}{अभिलक्षणिक
सदिश} $x \in G$ है (\cref{thm:b2:reduction:trigonalization}, अथवा सीधे किसी मूल का अस्तित्व); परंतु $u$ का हर
\omterm{def:b2:reduction:eigen}{अभिलक्षणिक सदिश} $F$ में पड़ता है, अतः $x \in F
\cap G = \{0\}$: विरोधाभास। अतः $G = \{0\}$ और
$E = F$: अभिलक्षणिक समष्टियाँ $E$ को भर देती हैं, अर्थात् $u$
\omterm{def:b2:reduction:diag}{विकर्णनीय} है।
\end{solution}

\begin{solution}{exo:b2:reduction:11}
त्रिभुजन कीजिए: $A = PTP^{-1}$, $T = \begin{pmatrix} \lambda & c\\ 0 &
\mu\end{pmatrix}$, $\abs\lambda, \abs\mu < 1$। तब $A^k =
PT^kP^{-1}$, और इतना ही पर्याप्त है कि $T^k \to 0$।

\emph{भिन्न \omterm{def:b2:reduction:eigen}{अभिलक्षणिक मान}:} आगमन से
\[
T^k = \begin{pmatrix}
\lambda^k & c\,\dfrac{\lambda^k - \mu^k}{\lambda - \mu}\\[4pt]
0 & \mu^k
\end{pmatrix},
\]
मिलता है, और हर प्रविष्टि $0$ की ओर जाती है ($\abs{\lambda}^k, \abs\mu^k \to 0$)।

\emph{बराबर \omterm{def:b2:reduction:eigen}{अभिलक्षणिक मान} ($\mu = \lambda$):} $T = \lambda I + cE_{12}$ और $T^k = \lambda^k I + k\lambda^{k-1}cE_{12}$; और प्रविष्टि $k\lambda^{k-1} \to 0$
क्योंकि $\abs\lambda < 1$ (गुणोत्तर बहुपद को हरा देता है)। दोनों स्थितियों में
प्रविष्टि-दर-प्रविष्टि $T^k \to 0$, अतः $A^k =
PT^kP^{-1} \to 0$ (स्थिर $P, P^{-1}$ से आव्यूह गुणन
प्रविष्टियों में संतत है --- गुणनफल की हर प्रविष्टि एक निश्चित रैखिक
संचय है)।
\end{solution}

\begin{solution}{exo:b2:reduction:12}
$\ker u$ की विमा $n - 1$ है (कोटि--शून्यता), अतः $0$ \omterm{def:b2:reduction:charpoly}{ज्यामितीय बहुलता} $n - 1$ वाला
\omterm{def:b2:reduction:eigen}{अभिलक्षणिक मान} है, और $\chi_u$ $X^{n-1}$ से विभाज्य है (\cref{def:b2:reduction:charpoly}: ज्यामितीय $\leq$
बीजीय)। $\chi_u = X^{n-1}(X - \alpha)$ लिखिए; $X^{n-1}$ का गुणांक $-\operatorname{tr} u$ होने से $\alpha =
\operatorname{tr} u$ मिलता है: $\chi_u = X^{n-1}(X - \operatorname{tr}
u)$।

यदि $\operatorname{tr} u \neq 0$: तो \omterm{def:b2:reduction:eigen}{अभिलक्षणिक मान} $\operatorname{tr} u$ $\chi_u$ का मूल है, अतः वह कोई \omterm{def:b2:reduction:eigen}{अभिलक्षणिक
सदिश} वहन करता है; $0$ और $\operatorname{tr} u$ की अभिलक्षणिक समष्टियों की विमाएँ $n - 1$
तथा $\geq 1$ हैं, जिनका योग $\geq n$ है: अतः वे $E$ को भर देती हैं और $u$
\omterm{def:b2:reduction:diag}{विकर्णनीय} है (\cref{thm:b2:reduction:diagcrit})। और यदि $\operatorname{tr} u = 0$: तो \cref{exo:b2:linalg:5} से $u \neq 0$ सहित $u^2 = (\operatorname{tr} u)u = 0$ मिलता है:
$u$ अशून्य शून्यंभावी है, और \omterm{def:b2:reduction:diag}{विकर्णनीय} शून्यंभावी शून्य होता है
(\cref{prop:b2:reduction:nilpotent}): अतः \omterm{def:b2:reduction:diag}{विकर्णनीय} नहीं।
\end{solution}

\begin{solution}{pb:b2:reduction:1}
\textbf{1.} $Cv_n$ के पहले $k - 1$ निर्देशांक $u_{n+1},
\dots, u_{n+k-1}$ हैं (अधिविकर्ण
स्थानांतरण कर देता है), और अंतिम $a_0 u_n + \dots + a_{k-1}u_{n+k-1}$ है। अतः $v_{n+1} = Cv_n$ सभी $n$ के लिए तभी
सत्य है जब अंतिम निर्देशांक सभी $n$ के लिए मेल खाएँ, अर्थात् तभी जब
$(\mathcal R)$ सत्य हो। इसे दोहराने पर $v_n = C^nv_0$।

\textbf{2.} $D_k(X) = \det(XI_k - C)$ को पहले स्तंभ के अनुदिश खोलिए: दो अशून्य प्रविष्टियाँ
$X$ (स्थान $(1,1)$) और $-a_0$ (स्थान $(k,1)$) हैं। पहला उपसारणिक गुणांकों $a_1, \dots, a_{k-1}$ के
लिए $D_{k-1}$ के आकार का है; दूसरा उपसारणिक विकर्ण $-1$ वाला ऊपरी
त्रिभुजाकार है: \omterm{def:b2:linalg:det}{सारणिक} $(-1)^{k-1}$, और स्थान से चिह्न $(-1)^{k+1}$। $k$ पर आगमन (आधार
$k = 1$: $X - a_0$) देता है
\[
D_k(X) = X\bigl(X^{k-1} - a_{k-1}X^{k-2} - \dots - a_1\bigr) -
a_0 = P(X).
\]
$\mu_C$ के लिए: चूँकि किसी भी बहुपद के लिए $Q(C^{\mathsf T}) = Q(C)^{\mathsf T}$, अतः $C$ और $C^{\mathsf T}$ के
\omterm{def:b2:linalg:annihilator}{विलोपक} एक ही हैं, इसलिए \omterm{def:b2:reduction:polyu}{न्यूनतम बहुपद} भी एक ही। $C^{\mathsf T}$ के लिए: स्तंभ
$C^{\mathsf T}e_1 = e_2$, \dots, $C^{\mathsf
T}e_{k-1} = e_k$ पढ़े जाते हैं, अतः $(e_1, C^{\mathsf T}e_1, \dots, (C^{\mathsf
T})^{k-1}e_1)$ विहित आधार है: स्वतंत्र।
तब घात $< k$ के किसी बहुपद $Q
\neq 0$ के लिए $Q(C^{\mathsf T})e_1 \neq 0$ (वह आधार सदिशों का अतुच्छ
संचय है): $\deg\mu \geq
k$। और $\deg P = k$ सहित $\mu \mid \chi = P$ होने से: $\mu_C = P$।

\textbf{3.} $v = (1, \lambda, \dots, \lambda^{k-1})^{\mathsf
T}$ के लिए: $Cv$ की पंक्तियाँ $1$ से $k-1$ तक $\lambda, \lambda^2, \dots,
\lambda^{k-1}$ देती हैं,
अर्थात् $v$ की पहली $k - 1$ प्रविष्टियों का $\lambda$ गुना; और अंतिम पंक्ति $\sum_m a_m\lambda^m = \lambda^k -
P(\lambda) = \lambda^k = \lambda\cdot\lambda^{k-1}$
देती है। अतः $Cv =
\lambda v$। विलोमतः, $i < k$ के लिए समीकरण $(Cx)_i = \lambda x_i$ $x_{i+1} = \lambda x_i$ पढ़े जाते हैं:
अर्थात् कोई भी \omterm{def:b2:reduction:eigen}{अभिलक्षणिक सदिश} $v$ के अनुपाती है --- हर \omterm{def:b2:reduction:eigen}{अभिलक्षणिक
समष्टि} की विमा ठीक $1$ है। \omterm{def:b2:reduction:diag}{विकर्णनीय} तभी जब अभिलक्षणिक समष्टियों की
विमाओं का योग $k$ हो (\cref{thm:b2:reduction:diagcrit}) तभी जब $k$ भिन्न \omterm{def:b2:reduction:eigen}{अभिलक्षणिक मान} हों तभी
जब $P$ के $k$ भिन्न मूल हों (\omterm{def:b2:reduction:eigen}{अभिलक्षणिक मान} $\chi_C = P$ के मूल ही हैं)।

\textbf{4.} प्रत्येक $(\lambda_i^n)_n$ $(\mathcal R)$ का हल है: $\lambda_i^{n+k} = \lambda_i^n\,\lambda_i^k =
\lambda_i^n\sum_m a_m\lambda_i^m$। स्वतंत्रता: $n = 0, \dots, k-1$ के लिए
लुप्त होने वाला संचय $\sum_i c_i\lambda_i^n = 0$ $c_i$ में एक वांडरमोंड निकाय (\cref{exo:b2:linalg:11}) है: अतः सभी
$c_i = 0$। हल-समष्टि की विमा $k$ है (प्रश्न 6, जिसकी उपपत्ति प्रारंभिक और
स्वतंत्र है): अतः $k$ स्वतंत्र हल आधार बनाते हैं, और निर्देशांक
अद्वितीय हैं।

\textbf{5.} $P = X^2 - X - 6 = (X - 3)(X + 2)$: सामान्य हल $u_n = A\,3^n + B(-2)^n$। आरंभिक शर्तें: $A + B = 1$, $3A -
2B = 8$: $A = 2$, $B = -1$:
\[
u_n = 2\cdot 3^n - (-2)^n .
\]
(जाँच: $u_2 = u_1 + 6u_0 = 14$ और $2\cdot9 - 4 = 14$।)

\textbf{6.} $P(S)\bigl((u_n)\bigr)$ अनुक्रम $n \mapsto
u_{n+k} - a_{k-1}u_{n+k-1} - \dots - a_0u_n$ है: वह तभी लुप्त होता है जब $(\mathcal R)$ सत्य हो,
अतः हल-समुच्चय $\ker P(S)$ है, जो एक उपसमष्टि है। प्रतिचित्रण $\ker P(S) \to \C^k$, $u \mapsto (u_0, \dots,
u_{k-1})$,
रैखिक है, एकैकी है (पुनरावृत्ति आगमन से पहले $k$ मानों से $u_{k}, u_{k+1}, \dots$ निर्धारित
कर देती है) और आच्छादक भी (किसी भी आरंभिक आँकड़े से $u_n$ को पुनरावर्ती
ढंग से परिभाषित कीजिए): अतः विमा $k$।

\textbf{7.} \cref{thm:b2:reduction:kernels} की उपपत्ति केवल इनका उपयोग करती है: $\C[X]$ में बेज़ू
सर्वसमिका, और यह तथ्य कि किसी स्थिर अंतःरूपांतरण के बहुपद परस्पर
क्रमविनिमेय होते हैं। इनमें से कोई भी परिवेशी समष्टि की विमा का उल्लेख
नहीं करता: अतः प्रमेयिका $S \in \mathcal{L}(\mathcal{S})$ के लिए अक्षरशः सत्य है। इसलिए
\[
\ker P(S) = \bigoplus_{i=1}^{r}
\ker\,(S - \lambda_i\,\mathrm{id})^{m_i}.
\]

\textbf{8.} $Q \in \C[X]$ के लिए: $(S -
\lambda)\bigl(Q(n)\lambda^n\bigr)_n$ का $n$-वाँ पद $Q(n{+}1)\lambda^{n+1} - \lambda Q(n)\lambda^n =
\lambda^{n+1}(\Delta Q)(n)$ है, जिसमें $\Delta Q = Q(X{+}1) - Q(X)$ की घात
$\deg Q - 1$ है (अग्र पद कट जाते हैं)। इसे दोहराने पर $(S - \lambda)^m\bigl(Q(n)\lambda^n\bigr) =
\bigl(\lambda^{n+m}(\Delta^m Q)(n)\bigr)_n$, और $\deg Q \leq m - 1$ होने पर $\Delta^m Q = 0$:
अर्थात् दाहिना समुच्चय केंद्रक में समाहित है। वह विमा $m$ की उपसमष्टि
है: अनुक्रम $(n^j\lambda^n)_n$, $0 \leq j < m$, स्वतंत्र हैं, क्योंकि सभी $n$ के लिए $\sum_j c_j n^j\lambda^n = 0$
होने से ($\lambda^n \neq 0$ से भाग देकर) बहुपद $\sum_j c_jX^j$ को हर $n \in \N$ पर लुप्त होना पड़ता
है, अतः शून्य होना पड़ता है। विलोमतः $\dim\ker(S -
\lambda)^m \leq m$: $(S - \lambda)^m = \sum_j
\binom mj(-\lambda)^{m-j}S^j$ खोलने पर समीकरण $(S - \lambda)^m u =
0$
$m$ कोटि की रैखिक पुनरावृत्ति है (अग्र गुणांक $1$), अतः प्रश्न 6 की
तरह $u$ $u_0, \dots, u_{m-1}$ से निर्धारित होता है। विमाओं की समता से बात पूरी हो जाती
है।

\textbf{9.} प्रश्न 7 और 8 मिलाइए: हर हल $\ker(S -
\lambda_i)^{m_i}$ के अवयवों के योग के रूप
में अद्वितीय रूप से अपघटित होता है, अर्थात् $\deg Q_i \leq m_i - 1$ सहित $u_n = \sum_i Q_i(n)\lambda_i^n$; और $Q_i$
अद्वितीय हैं क्योंकि अपघटन प्रत्यक्ष है और प्रत्येक पद के भीतर $Q_i$ के
गुणांक आधार $(n^j\lambda_i^n)_j$ में निर्देशांक हैं (प्रश्न 8)। विमाओं की जाँच: $\sum_i m_i = k$।

\textbf{10.} $P = X^2 - 4X + 4 = (X - 2)^2$: हल $(a +
bn)2^n$। आरंभिक आँकड़े: $a = 1$, $2(a + b) = 0$, अतः $b = -1$:
\[
u_n = (1 - n)\,2^n .
\]
जाँच: $u_2 = 4u_1 - 4u_0 = -4$, और $(1 - 2)\cdot4 = -4$।

\textbf{11.} $u_n = \lambda_1^n\bigl(c_1 + \sum_{i\geq2}
c_i(\lambda_i/\lambda_1)^n\bigr)$ लिखिए; हर अनुपात का मापांक $< 1$ है, अतः कोष्ठक $c_1 \neq 0$ की
ओर जाता है: $u_n \sim c_1\lambda_1^n$। विशेष रूप से बड़े $n$ के लिए $u_n \neq 0$, और
\[
\frac{u_{n+1}}{u_n} =
\lambda_1\,\frac{c_1 + o(1)}{c_1 + o(1)} \longrightarrow
\lambda_1 .
\]

\textbf{12.} परिकलन कीजिए:
\[
q(a_{n+1}, b_{n+1}) = (a_n + 2b_n)^2 - 2(a_n + b_n)^2
= -a_n^2 + 2b_n^2 = -q(a_n, b_n).
\]
$q(a_0, b_0) = 1 - 2 = -1$ के साथ: $a_n^2 - 2b_n^2 = (-1)^{n+1}$। संरचनात्मक रूप से: $q(a, b) = (a - \sqrt2\,b)(a + \sqrt2\,b)$, और रैखिक प्रतिचित्रण $M$
गुणनखंड $a + \sqrt2 b$ को $1 +
\sqrt2$ से तथा गुणनखंड $a - \sqrt2 b$ को $1 - \sqrt2$ से गुणा कर देता है
(परिकलन: $a_{n+1} + \sqrt2 b_{n+1} = (1 + \sqrt2)(a_n + \sqrt2 b_n)$); अतः हर पग पर गुणनफल $(1 + \sqrt2)(1 - \sqrt2) = -1 =
\det M$ से गुणा हो जाता है।

\textbf{13.} चूँकि $a_n^2 - 2b_n^2 = (a_n - \sqrt2 b_n)(a_n +
\sqrt2 b_n) = (-1)^{n+1}$,
\[
\Bigl|\frac{a_n}{b_n} - \sqrt2\Bigr|
= \frac{\abs{a_n^2 - 2b_n^2}}{b_n(a_n + \sqrt2 b_n)}
= \frac{1}{b_n(a_n + \sqrt2 b_n)} \leq \frac1{2b_n^2},
\]
जिसमें $a_n \geq b_n \geq 1$ का उपयोग हुआ (आगमन: दोनों बढ़ते हैं), अतः $a_n +
\sqrt2 b_n \geq (1 + \sqrt2)b_n \geq 2b_n$। $M$ के
\omterm{def:b2:reduction:eigen}{अभिलक्षणिक मान} $\chi_M = X^2 - 2X - 1$ हैं, जिनके मूल $1 \pm \sqrt2$; और चूँकि $(a_0, b_0)$ का प्रबल
\omterm{def:b2:reduction:eigen}{अभिलक्षणिक सदिश} पर अशून्य घटक है (सभी प्रविष्टियाँ धनात्मक हैं), अतः
$c > 0$ सहित $b_n \sim c(1 + \sqrt2)^n$ (प्रश्न 11)। इसलिए त्रुटि $\asymp (1 + \sqrt2)^{-2n} = (3 +
2\sqrt2)^{-n}$ है: अर्थात् अनुपात $1/(3 + 2\sqrt2) = 3
- 2\sqrt2 \approx 0.172$
के साथ गुणोत्तर क्षय।

\textbf{14.} (क) प्रश्न 9 से: $\abs{u_n} \leq \sum_i
\abs{Q_i(n)}\abs{\lambda_i}^n \leq \bigl(\sum_i
\abs{Q_i(n)}\bigr)\rho^n$, और $n \geq 1$ के लिए हर $\abs{Q_i(n)} \leq C_i
n^{m_i - 1} \leq C_i n^{m-1}$: अचरों को जोड़
दीजिए। (ख) मान लीजिए $\rho' = \max_{i \geq 2}\abs{\lambda_i} <
\abs{\lambda_1}$ और $d = \deg Q_1$, जिसका अग्र गुणांक $c
\neq 0$ है। तब $\abs{R_n}
\leq Cn^{m-1}\rho'^n$
सहित $u_n = Q_1(n)\lambda_1^n + R_n$, और
\[
\frac{R_n}{Q_1(n)\lambda_1^n} = O\Bigl(n^{m-1-d}
\bigl(\rho'/\abs{\lambda_1}\bigr)^n\Bigr) \longrightarrow 0
\]
(गुणोत्तर बहुपद को हरा देता है)। अतः
\[
u_n \sim Q_1(n)\,\lambda_1^n \sim c\,n^d\lambda_1^n,
\qquad
\frac{u_{n+1}}{u_n} \longrightarrow \lambda_1
\quad\text{(क्योंकि } Q_1(n{+}1)/Q_1(n) \to 1\text{)}.
\]
प्रश्न 10 पर जाँच: $u_n = (1-n)2^n$ के लिए अनुपात
\[
\frac{(-n)2^{n+1}}{(1-n)2^n} = 2\,\frac{-n}{1-n}
\longrightarrow 2 = \lambda_1 .
\]

\textbf{15.} $n$ पर आगमन। $n = 1$ के लिए $A_{ij}$ $1$ लंबाई के पथ गिनता है।
चरण: $i$ से $j$ तक $n + 1$ लंबाई का पथ $i$ से किसी शीर्ष $\ell$ तक $n$
लंबाई का पथ है जिसके बाद एक कोर $\ell j$ आती है:
\[
\#\{\text{पथ}\} = \sum_{\ell} (A^n)_{i\ell}A_{\ell j} =
(A^{n+1})_{ij}.
\]

\textbf{16.} $J = 3\Pi$, जहाँ $\Pi = J/3$ समतल $x + y + z = 0$ के अनुदिश $\operatorname{Vect}(1,1,1)$ पर प्रक्षेप है
($\Pi^2 = \Pi$ क्योंकि $J^2 = 3J$)। तब $A = J - I = 2\Pi - (I -
\Pi)$, और चूँकि $\Pi$ तथा $I - \Pi$ पूरक प्रक्षेप हैं,
\[
A^n = 2^n\,\Pi + (-1)^n (I - \Pi),
\qquad\text{अर्थात्}\qquad
(A^n)_{ij} = \frac{2^n}3 + (-1)^n\Bigl(\delta_{ij} -
\frac13\Bigr),
\]
जिससे दोनों प्रदर्शित सूत्र मिल जाते हैं। $n = 2$ पर: विकर्ण $(4 + 2)/3 = 2$ (दो
पड़ोसियों $\ell$ के लिए पथ $i \to \ell \to i$) और विकर्णेतर $(4 - 1)/3 = 1$ (तीसरे शीर्ष से होकर
जाने वाला अकेला पथ $i \to
\ell \to j$)।

\textbf{17.} मान लीजिए $w_n^{(0)}, w_n^{(1)}$ $n$ लंबाई के उन ग्राह्य शब्दों को गिनते हैं
जो क्रमशः $0$ और $1$ पर समाप्त होते हैं। एक अक्षर जोड़ने पर: $0$ किसी
के भी बाद आ सकता है, जबकि $1$ केवल $0$ के बाद:
\[
\begin{pmatrix} w_{n+1}^{(0)}\\ w_{n+1}^{(1)}\end{pmatrix}
= \begin{pmatrix} 1 & 1\\ 1 & 0\end{pmatrix}
\begin{pmatrix} w_n^{(0)}\\ w_n^{(1)}\end{pmatrix}.
\]
जोड़ने पर $w_{n+2} = w_{n+1} + w_n$ (अथवा: पहले अक्षर पर शर्त लगाइए)। $w_1 = 2$, $w_2 = 3$ के साथ:
आगमन से $w_n = F_{n+2}$ ($F_3 = 2$, $F_4 = 3$, और वही पुनरावृत्ति)। वृद्धि: $X^2 - X - 1$ के मूल
$\varphi > \abs\psi$ हैं (\cref{exo:b2:reduction:5}), और $\varphi$-घटक अशून्य है ($w_n$ धनात्मक हैं और $\psi^n \to 0$),
अतः प्रश्न 11 से $w_{n+1}/w_n \to \varphi = \frac{1 + \sqrt5}2$।

\textbf{18.} $A = \begin{psmallmatrix} 0&1&0\\ 1&0&1\\
0&1&0\end{psmallmatrix}$। जाँच:
\[
A(1, \pm\sqrt2, 1)^{\mathsf T}
= (\pm\sqrt2, 2, \pm\sqrt2)^{\mathsf T}
= \pm\sqrt2\,(1, \pm\sqrt2, 1)^{\mathsf T},
\qquad
A(1, 0, -1)^{\mathsf T} = 0 :
\]
\omterm{def:b2:reduction:eigen}{अभिलक्षणिक मान} $\sqrt2, -\sqrt2, 0$ ($= 2\cos\frac\pi4,
2\cos\frac{3\pi}4, 2\cos\frac\pi2$)। $e_1$ को अभिलक्षणिक आधार पर अपघटित करके तीसरा
निर्देशांक पढ़िए, अथवा सममिति का उपयोग कीजिए: $v_\pm = (1, \pm\sqrt2, 1)$, $v_0 = (1, 0, -1)$ के साथ $e_1
= \frac14 v_+ + \frac14 v_- + \frac12 v_0$
जाँचा जा सकता है, अतः $n \geq 1$ के लिए
\[
(A^n)_{13} = \Bigl(\tfrac14(\sqrt2)^n v_+ +
\tfrac14(-\sqrt2)^n v_- + 0\Bigr)_{\!3}
= \frac{(\sqrt2)^n + (-\sqrt2)^n}{4},
\]
जो विषम $n$ के लिए शून्य है (द्विभाजित ग्राफ: दोनों सिरे सम दूरी पर
हैं), और सम $n$ के लिए $2\cdot 2^{n/2}/4 = 2^{n/2 - 1}$। $n = 4$ पर: $2^{1} = 2$, जो दोनों पथों $1\,2\,1\,2\,3$ और
$1\,2\,3\,2\,3$ से मेल खाता है।

\textbf{19.} $i$ से $n$ लंबाई के बंद पथ $(A^n)_{ii}$ हैं; $i$ पर जोड़ने से $\operatorname{tr}(A^n)$
मिलता है। $A$ का ($\C$ पर) त्रिभुजन करने पर $A^n$ विकर्ण $\lambda_i^n$ वाला
त्रिभुजाकार है: $\operatorname{tr}(A^n) = \sum_i\lambda_i^n$। त्रिभुज: $\operatorname{tr}(A^n) = 3\,\frac{2^n + 2(-1)^n}3 =
2^n + 2(-1)^n = 2^n + (-1)^n + (-1)^n$: \omterm{def:b2:reduction:eigen}{वर्णक्रम} $\{2, -1,
-1\}$, जो प्रश्न 16 के
अनुरूप है।

\textbf{20.} $W^{\mathsf T}$ के स्तंभ: $i \geq 1$ के लिए $W^{\mathsf T}e_i =
e_{i-1}$ और $W^{\mathsf T}e_0 = e_{k-1}$; क्रम $e_0, e_1, \dots$ में पुनः
अंकित करने पर यह ठीक $X^k - 1$ का सहचर आव्यूह है ($a_0 = 1$, शेष $a_m = 0$)। प्रश्न 2:
$\chi_{W} = \chi_{W^{\mathsf T}} = X^k - 1 =
\mu_{W}$। मूल $\omega^j$ ($j = 0, \dots, k-1$) इकाई के $k$ भिन्न $k$-वें मूल हैं: अतः $W$
\omterm{def:b2:reduction:diag}{विकर्णनीय} है (प्रश्न 3, अथवा \cref{exo:b2:reduction:8}: $W^k = I$)। \omterm{def:b2:reduction:eigen}{अभिलक्षणिक सदिश}: $Wf_j
= \sum_m \omega^{-jm}e_{m+1} = \sum_{m'}\omega^{-j(m'-1)}e_{m'}
= \omega^j f_j$।

\textbf{21.} \omterm{def:b2:structures:generated}{चक्रीय} आव्यूह $W$ के बहुपद हैं, और किसी स्थिर आव्यूह के
बहुपद परस्पर क्रमविनिमेय होते हैं। प्रत्येक $f_j$ हर घात का \omterm{def:b2:reduction:eigen}{अभिलक्षणिक
सदिश} है: $W^m f_j = \omega^{jm}f_j$, अतः
\[
Cf_j = \sum_m c_m\omega^{jm} f_j = \widehat c(\omega^j)\,f_j :
\]
और आधार $(f_0, \dots, f_{k-1})$ (स्वतंत्र: भिन्न $\omega^{-j}$, \cref{exo:b2:linalg:11} का वांडरमोंड) हर \omterm{def:b2:structures:generated}{चक्रीय} आव्यूह
का एक ही बार में विकर्णन कर देता है, और \omterm{def:b2:reduction:eigen}{अभिलक्षणिक मान} कथनानुसार होते
हैं।

\textbf{22.} \omterm{def:b2:linalg:det}{सारणिक} अभिलक्षणिक मानों का गुणनफल है (विकर्णन कीजिए):
$\det C = \prod_{j}\widehat c(\omega^j)$। $k
= 3$ के लिए $c_0 = a$, $c_1 = b$, $c_2 = c$ और $\omega = j =
\eu^{2\iu\pi/3}$:
\[
\det C = (a + b + c)(a + bj + cj^2)(a + bj^2 + cj^4),
\]
तथा $j^4 = j$: अर्थात् ठीक \cref{exo:b2:linalg:8} का गुणनखंडन।

\textbf{23.} $M = \frac12(W + W^{-1})$ एक \omterm{def:b2:structures:generated}{चक्रीय} आव्यूह है ($W^{-1} =
W^{k-1}$), जिसके \omterm{def:b2:reduction:eigen}{अभिलक्षणिक मान} उसी
आधार $f_j$ पर $\frac12(\omega^j + \omega^{-j}) =
\cos\frac{2\pi j}k$ हैं। निर्देशांक: $x^{(0)} = \sum_j \alpha_j f_j$ लिखिए। $f_j$ के निर्देशांकों का योग
$\sum_m \omega^{-jm}$ है, जो $j = 0$ के लिए $k$ और अन्यथा $0$ होता है (अनुपात $\omega^{-j} \neq 1$ वाला
गुणोत्तर योग)। $x^{(0)}$ के निर्देशांक जोड़ने पर: $\sum_m x^{(0)}_m =
\alpha_0\,k$, अतः $\alpha_0 = \frac1k\sum_m x^{(0)}_m$, अर्थात्
माध्य।

\textbf{24.} $x^{(n)} = M^nx^{(0)} = \sum_j
\alpha_j\cos^n\bigl(\tfrac{2\pi j}k\bigr)f_j$। विषम $k$ के लिए हर $j \neq 0$ पर $\abs{\cos(2\pi j/k)} < 1$ (कोण कभी $0$ या $\pi$
नहीं होता), अतः $j = 0$ को छोड़कर सभी पद $0$ की ओर जाते हैं: $x^{(n)} \to \alpha_0 f_0$, अर्थात्
वह अचर सदिश जो माध्य के बराबर है --- विषम वलय पर औसत लेना सबको बराबर
कर देता है। $k = 4$ के लिए \omterm{def:b2:reduction:eigen}{अभिलक्षणिक मान} $1, 0, -1, 0$ हैं: $f_2 = (1, -1, 1, -1)^{\mathsf T}$ के साथ $j = 2$ वाला
पद $\alpha_2(-1)^nf_2$ सदा दोलन करता रहता है। बाधा एकांतर माध्य है: $x^{(0)}$ के
निर्देशांकों को $(-1)^m$ से गुणा करके जोड़ने पर वही गुणोत्तर-योग परिकलन
$\sum_m
(-1)^mx^{(0)}_m = 4\alpha_2$ देता है: अतः प्रक्रिया अभिसरित होती है यदि और केवल यदि $x^{(0)}_0 - x^{(0)}_1 + x^{(0)}_2 - x^{(0)}_3 = 0$, और
तब वह माध्य पर अभिसरित होती है।

\textbf{25.} सहचर आव्यूह कोटि $k$ की अदिश पुनरावृत्ति को प्रथम कोटि की
सदिश पुनरावृत्ति में बदल देता है, जिससे बंद सूत्र $C^n$ के बारे में कथन
बन जाते हैं --- और यह समानयन का अपना मैदान है (प्रश्न 1--5)। केंद्रक
अपघटन प्रमेयिका शुद्ध बहुपद बीजगणित है (बेज़ू और क्रमविनिमेयता), इसलिए
वह $\ker P(S)$ को तब भी बाँट देती है जब $\mathcal{S}$ अनंत-विमीय हो (प्रश्न 7--9)। प्रबल
\omterm{def:b2:reduction:eigen}{अभिलक्षणिक मान} वृद्धि पर शासन इसलिए करते हैं कि सामान्यीकरण के बाद हर
दूसरा योगदान गुणोत्तर रूप से नगण्य हो जाता है --- और इसी से पेल की
त्रुटि प्रबल मूल के वर्ग की दर से क्षीण होती है (प्रश्न 11--14)।
संलग्नता आव्यूह की घातें पथ इसलिए गिनती हैं कि आव्यूह गुणन मध्यवर्ती
शीर्षों पर योग ले लेता है, इसलिए \omterm{def:b2:reduction:eigen}{वर्णक्रम} बंद पथ गिन देते हैं (प्रश्न
15--19)। क्रमविनिमेय आव्यूहों का अभिलक्षणिक आधार साझा होता है, और तब एक
ही फूरिये आधार पूरे \omterm{def:b2:structures:generated}{चक्रीय} बीजगणित का एक ही झटके में विकर्णन कर देता है
(प्रश्न 20--24)। शिखर: रैखिक पुनरावृत्तियों की मूल प्रमेय (प्रश्न 9);
और अऋणात्मक आव्यूहों के लिए, $\varphi$ या $1 +
\sqrt2$ जैसे प्रबल मूल स्वतः वास्तविक,
धनात्मक और सरल क्यों होते हैं इसका कारण पेरों--फ्रोबेनियस प्रमेय है,
जिसे तृतीय वर्ष के खंड में सिद्ध किया गया है।
\end{solution}
